\documentclass[11pt]{article}

\usepackage[margin=1in]{geometry}
\usepackage{amsmath,amssymb,amsthm}
\usepackage{mathtools}
\usepackage{enumitem}
\usepackage{hyperref}

\newtheorem{theorem}{Theorem}
\newtheorem{lemma}[theorem]{Lemma}
\newtheorem{proposition}[theorem]{Proposition}
\newtheorem{corollary}[theorem]{Corollary}
\theoremstyle{definition}
\newtheorem{definition}[theorem]{Definition}
\newtheorem{axiom}{Axiom}
\renewcommand{\theaxiom}{M\arabic{axiom}}
\theoremstyle{remark}
\newtheorem{remark}[theorem]{Remark}

\newcommand{\OK}{\mathcal{O}_K}
\newcommand{\Odelta}{\mathcal{O}_\delta}
\newcommand{\Fq}{\mathbb{F}_q}
\newcommand{\Fr}{\mathbb{F}_r}
\newcommand{\Zp}{\mathbb{Z}_p}
\newcommand{\Qp}{\mathbb{Q}_p}
\newcommand{\Monic}{\mathrm{Monic}}
\newcommand{\type}{\mathrm{type}}
\newcommand{\disc}{\mathrm{disc}}
\newcommand{\Res}{\mathrm{Res}}
\newcommand{\NP}{\mathrm{NP}}
\newcommand{\rhocount}{\rho}
\newcommand{\vp}{v}
\DeclareMathOperator{\dvg}{\mathbf{v}}

\title{A counting-model blueprint for the $p$-adic density decomposition:\\
density as a sum over the OM tree, with the geometric collapse \\
\large (derived from one Montes axiom $+$ elementary counting)}
\author{Blueprint for hand-verification before formalization}
\date{2026-06-19}

\begin{document}
\maketitle

\begin{abstract}
We give a precise, self-contained derivation that the $p$-adic factorization-type
density $\rhocount(n,\sigma;q)$ is a \emph{finite sum over the combinatorial
Okutsu--Montes (OM) tree-shapes} of rational coefficients $C_T(q)$, \emph{and}
that each $C_T(q)$ is obtained from an \emph{infinite} literal OM tree by a
\emph{self-loop collapse}. The collapse is a single geometric series
$V=a+rV\Rightarrow V=a/(1-r)$ \emph{only} at a size-$2$ self-loop; at a size-$s\ge3$
self-loop it is a \emph{well-founded nested system} of geometric series (a
size-$s$ self-loop pivot $1/(1-q^{-w(s)})$ applied to a finite sum over
$\ge2$-block partitions whose blocks carry \emph{strictly smaller} child volumes),
and the slowest mode of the depth-graded tail is $q^{-2}$, \emph{not} $q^{-w(s)}$.
Each $C_T(q)$ is nonetheless rational with denominator a product of cyclotomic
factors $(q^{w(s)}-1)$. The single imported assumption is a
\emph{count-multiplicativity} axiom (axiom~\ref{ax:montes}, ``what Montes says
about factorization given the Newton polygon''); everything else --- finiteness
of shapes, the partition, counting-additivity, the nested collapse, and
rationality --- is \emph{proved} from that axiom plus elementary counting.
The whole argument is run in a \textbf{counting model} (\S\ref{sec:counting}) in
which the abstract $p$-adic Haar ``measure wall'' never appears: additivity is
counting-additivity, ball/box volumes are residue-class counts, and the
collapse is a finite geometric sum (resp.\ a finite nested system) followed by an
honest $N\to\infty$ limit. Every closed-form claim below is confirmed by an
\emph{exact} (rational-arithmetic, never Monte-Carlo) enumerator over $\Zp/p^N$;
see \S\ref{sec:formalization}, \texttt{notes/decomp\_blueprint\_verify.py} and
\texttt{notes/r1\_verify.py}.
\end{abstract}

\paragraph{Changelog (2026-06-19, after adversarial exact-enumeration review).}
The crux was restated to fix a structural error and several normalization slips
surfaced by independent exact (rational-arithmetic) enumeration over $(\Zp/p^N)^s$
at $p\in\{2,3,5,7\}$:
\begin{enumerate}[label=\textup{[F\arabic*]},nosep,leftmargin=*]
\item\textbf{The size-$s$ self-loop is NOT a single geometric series for $s\ge3$.}
The previous draft wrote $V=a+rV$ with a \emph{constant} $a$ and a \emph{single}
ratio $r=q^{-w(s)}$ at every $s$. Exact enumeration shows this is true only for
$s=2$. For $s=3$ the lumped $a_N$ \emph{strictly increases} with precision $N$
(e.g.\ $p{=}2$: $L_N=0,\tfrac3{16},\tfrac{123}{512},\dots$, not constant) and the
depth-increment ratio $D_N/D_{N-1}\to q^{-2}=\tfrac14$ (the nested size-$2$
sub-cluster), \emph{not} $q^{-w(3)}=q^{-5}=\tfrac1{32}$. The crux is now stated as
the \emph{nested multi-block recursion} (Theorem~\ref{thm:nested}); the single
geometric series is its $s=2$ special case (Lemma~\ref{lem:s2}).
\item\textbf{The verifier's $s=3$ block now matches.} The old script ran a
single-geometric extrapolator on $s=3$ and \emph{failed}; the new
\texttt{decomp\_blueprint\_verify.py} (i)~confirms $s=2$ exactly, (ii)~exhibits
the $s\ge3$ failure of the single-geometric model as a deliberate demonstration,
and (iii)~checks the nested recursion against the closed forms and against a
rigorous exact bracket $[L_N,L_N+U_N]$: $I_3=8/31\in[0.25800,0.28131]$ at $p{=}2$,
$189/484\in[0.39026,0.42699]$ at $p{=}3$.
\item\textbf{Corrected the boxed $s=2$ coefficient:} $a=(q-1)/q$ exactly
(verified $a=\tfrac12,\tfrac23,\tfrac45$ at $p=2,3,5$), and $a/(1-q^{-2})=q/(q+1)$.
The garbled product on the old line~417 is removed.
\item\textbf{Normalization caveat on $\rhocount$:} the assembly
$(q^2+1)^2/(6\Phi_5)$ is the \emph{projective / $3$-distinct-points} density, not
the monic counting density of Definition~\ref{def:count}; it is now labeled as
such (it uses the symmetry factor $3!$ and division by $|\mathbb P^3|$). The monic
type densities are different rationals (e.g.\ inert quadratic
$=(q-1)/(2q)$, exactly verified).
\item\textbf{Squarefree/a.e.\ caveat} added to Definition~\ref{def:count}: the
type is defined off the measure-zero discriminant locus.
\end{enumerate}

\paragraph{Changelog (2026-06-21, box-wise refactor of the Montes input).}
The single Lean field that previously bundled the entire per-node count recursion
(\texttt{MontesData.nodeMultiplicativity}) has been \emph{decomposed} into box-wise
measure-theoretic primitives, three of which are now \emph{genuinely proved} in Lean
(core-only \texttt{\#print axioms}), leaving \emph{one} minimal irreducible
box-Haar/equidistribution assumption. \S\ref{sec:ledger} (new) records the resulting
\textbf{exact assumed-vs-proven ledger} item by item, matching the
\texttt{\#print axioms} reality of the current Lean development. In short: the
per-node factorization \eqref{eq:nodefactor} is now a \emph{Lean theorem}
(\texttt{nodeMultiplicativity}, core-only) derived from a proved box-volume
($T\_BB1$), a proved residual-config count ($T\_BB3$), proved box-additivity, and the
single minimal axiom; the abstract $p$-adic Haar identification is the only
irreducible residue, isolated to the per-box fields
\texttt{boxHaarEquidist}\,/\,\texttt{nodeMeasure\_boxSum}. The headline theorem
\texttt{goal\_theorem\_montes} carries exactly
$\{\texttt{propext},\ \texttt{Classical.choice},\ \texttt{Quot.sound},\
\texttt{tame\_functionalEquation}\}$; the rationality/decomposition constants carry
only the three Lean-core axioms. \S\ref{sec:ledger} also records the resolved
non-vacuity witness (the full coupled bundle is now sorry-free inhabited).

\tableofcontents

\section{Setup and the counting model}\label{sec:counting}

\subsection{Conventions}

Let $K/\Qp$ be a finite \emph{unramified} extension with residue field $\Fq$
($q=p^k$), ring of integers $\OK$, uniformizer $\pi$ (we may take $\pi=p$),
normalized valuation $\vp$ with $\vp(p)=1$, and Haar measure $\mu$ normalized by
$\mu(\OK)=1$. Unramifiedness is essential and is used throughout: $\vp(p)=1$ is
the one normalization constant that carries all the characteristic dependence, and
it is preserved under unramified base change. For the unramified extension
$K_\delta/K$ of degree $\delta$ we write $\Odelta$, residue field $\mathbb F_{q^\delta}$,
$Q:=q^\delta$, and use the index identity
\begin{equation}\label{eq:index}
  [\Odelta : p^m\Odelta] = Q^m \qquad\text{for every prime, the same $Q$.}
\end{equation}

A \emph{factorization type} $\sigma$ of degree $n$ is the multiset of pairs
$(e_i,f_i)$ ($\sum_i e_if_i=n$) of ramification/residue indices of the irreducible
$p$-adic factors of a monic degree-$n$ polynomial. The object of interest is the
\emph{density} of monic polynomials of type $\sigma$.

\subsection{The counting-model definition of density}

The classical definition would use Haar measure on the coefficient space
$\OK^n\cong\{f=x^n+a_{n-1}x^{n-1}+\dots+a_0\}$. mathlib has no $p$-adic Haar
measure (and constructing it is the ``measure wall''), and even informally the
additivity/ball-volume bookkeeping is where the abstract argument gets heavy. We
sidestep this by \emph{defining} the density as a limit of finite normalized counts.

\begin{definition}[Counting-model density]\label{def:count}
Fix $n$ and a type $\sigma$. For a precision $N\ge 1$, identify a monic
$f\in(\OK/p^N)[x]$ of degree $n$ with its coefficient tuple
$(a_0,\dots,a_{n-1})\in(\OK/p^N)^n$, so there are exactly $q^{nN}$ of them. Say
$f$ \emph{has type $\sigma$ resolved at precision $N$} if every lift of $f$ to a
monic polynomial over $\OK$ has factorization type $\sigma$ --- i.e.\ the type is
\emph{decided} by the digits below $p^N$. Define
\begin{equation}\label{eq:rhodef}
  \rhocount(n,\sigma;q)\;:=\;\lim_{N\to\infty}\,
  q^{-nN}\cdot\#\bigl\{\,\text{monic }f\in(\OK/p^N)[x]\text{ of degree }n
  : \type(f)=\sigma\text{ resolved at }N\,\bigr\}.
\end{equation}
The factorization type is only well-defined on the \emph{separable} (squarefree)
locus $\{\disc f\ne0\}$; on the complementary $\{\disc f=0\}$ locus (repeated
irreducible factors) the type as a multiset of $(e_i,f_i)$ is not determined by
finitely many digits. That locus is measure-zero (it is cut out by the nonzero
polynomial $\disc f$), so it is irrelevant to the limit \eqref{eq:rhodef}; the
``resolved at $N$'' clause automatically excludes it (a polynomial with a repeated
factor is never decided), and it contributes to $U_N\to0$ rather than to any
$L_N$. We always work on the squarefree locus.
\end{definition}

\begin{lemma}[The counting model computes the Haar density]\label{lem:countHaar}
The limit \eqref{eq:rhodef} exists and equals the Haar volume
$\mu\{f\in\OK^n:\type(f)=\sigma\}$, whenever the latter is defined (i.e.\ off a
measure-zero locus).
\end{lemma}

\begin{proof}
Write $L_N := q^{-nN}\cdot\#\{f:\type=\sigma\text{ decided at }N\}$ and
$U_N := q^{-nN}\cdot\#\{f:\text{undecided at }N\}$. A coset
$f+p^N(\OK/\cdots)$ of measure $q^{-nN}$ that is \emph{decided} (its type is
determined by digits $<p^N$) has \emph{all} of its lifts of type $\sigma$ or all
not; the cylinder $\{g\in\OK^n: g\equiv f\ (p^N)\}$ has Haar measure exactly
$q^{-nN}$ by \eqref{eq:index} (a product of $n$ one-coordinate balls). Hence $L_N$
is a lower bound and $L_N+U_N$ an upper bound for the Haar volume:
\[
  L_N \;\le\; \mu\{\type=\sigma\}\;\le\; L_N+U_N .
\]
$L_N$ is nondecreasing (refining a decided coset keeps it decided), so
$\lim_N L_N$ exists. The undecided locus at precision $N$ is contained in the
\emph{boundary} where $f$ has a repeated factor not yet separated, i.e.\ a subset
of the discriminant-vanishing locus $\{\,\vp(\disc f)\ge N-c\,\}$ for a constant
$c=c(n)$; this has Haar measure $\to 0$ (the discriminant is a nonzero polynomial,
so $\{\vp(\disc f)\ge M\}$ has measure $O(q^{-\lfloor M/2\rfloor})\to0$). Thus
$U_N\to0$ and $L_N\to\mu\{\type=\sigma\}$.
\end{proof}

\begin{remark}[Why this sidesteps the measure wall]\label{rem:wall}
In the counting model every measure-theoretic primitive used below is an
\emph{elementary finite count}:
\begin{itemize}[nosep]
\item \textbf{Additivity} is counting-additivity over a finite partition of
$(\OK/p^N)^n$ --- trivial, no $\sigma$-algebra needed.
\item \textbf{Ball / box volume} is a residue-class count divided by $q^{nN}$:
$\mu\{\,\vp(c)\ge h\,\}=q^{-h}$ because the condition $c\equiv0\ (p^h)$ singles
out $q^{N-h}$ of the $q^N$ residues mod $p^N$ in that coordinate, and
$q^{N-h}/q^N=q^{-h}$; similarly $\mu\{\vp(c)=h\}=(1-q^{-1})q^{-h}$. These are
\eqref{eq:index} restated as counting.
\item \textbf{The geometric collapse} (\S\ref{sec:collapse}) is a \emph{finite}
geometric sum of such counts at each precision $N$, followed by the honest
$N\to\infty$ limit of Definition~\ref{def:count}.
\end{itemize}
We never invoke an abstract Haar measure; Lemma~\ref{lem:countHaar} is the only
bridge to the classical statement, and it is itself elementary.
\end{remark}

\section{The Montes axiom (the one imported assumption)}\label{sec:axiom}

We import \emph{exactly one} fact, the count-shadow of Montes' factorization
theorem ``given the Newton polygon''. We state it as an axiom and mark it (A).
Everything in \S\S\ref{sec:decomp}--\ref{sec:rational} is then (B), proved from
(A) plus elementary counting.

\subsection{The OM tree of strata}

\begin{definition}[OM tree and shapes]\label{def:omtree}
Following Gu\`ardia--Montes--Nart (GMN, \emph{Newton polygons of higher order in
algebraic number theory}, Trans.\ Amer.\ Math.\ Soc.\ \textbf{364} (2012)
361--416; arXiv:0807.2620), a monic $f\in\OK[x]$ carries OM/Montes data organized
as a rooted tree. Each \emph{node} of order $r$ records:
\begin{enumerate}[label=(\roman*),nosep]
\item an order-$r$ \emph{key polynomial} $\phi_r$ (a monic lift of the previous
residual factor);
\item a \emph{side} of the order-$r$ $\phi_r$-Newton polygon $N_r$, of slope
$-h/e$ in lowest terms (a lattice polygon in $\mathbb Z\times\frac1{D_r}\mathbb Z$,
$D_r=e_1\cdots e_{r-1}$ a $p$-independent integer);
\item a \emph{residual polynomial} $R_r\in\Fr[y]$ over the finite order-$r$
residue field $\Fr=\mathbb F_{q^{\delta f_1\cdots f_{r-1}}}$, together with its
\emph{factorization pattern} (degrees and multiplicities of factors).
\end{enumerate}
\emph{Children} of a node are the descents on its \emph{repeated} residual
factors (multiplicity $\ge2$); a \emph{leaf} is a multiplicity-$1$ residual
factor, an irreducible $p$-adic factor with explicit data $(e^\ast,f^\ast)$.
The \emph{combinatorial shape} $T$ records only this decorated tree --- orders,
slopes $(h,e)$, residual factorization \emph{patterns} (degrees/multiplicities,
not the residual polynomials themselves), and residue degrees $f_i$ --- not the
specific coefficients. We write $\type(T)$ for the type $\sigma$ assembled from
the leaf data $(e^\ast_j,f^\ast_j)$.
\end{definition}

For a fixed shape $T$ let $S_T^{(N)}\subseteq(\OK/p^N)^n$ be the set of monic $f$
whose OM tree (resolved at precision $N$) realizes $T$; the \emph{stratum} of $T$.

\subsection{The axiom}

\begin{axiom}[Montes count-multiplicativity over the OM tree]\label{ax:montes}
The OM data determines the factorization, and the count of monic polynomials lying
in a tree-stratum \emph{factors over the tree}. Precisely:
\begin{enumerate}[label=\textup{(\arabic*)},nosep]
\item\label{ax:partition} \textbf{(Unique tree / partition.)} Every monic $f$ has a
unique OM tree, so for each $N$ the strata $\{S_T^{(N)}\}_T$ (over combinatorial
shapes $T$, together with the residual-pattern decorations resolved at $N$)
\emph{partition} the monic polynomials of degree $n$ in $(\OK/p^N)[x]$.
\item\label{ax:multiplicativity} \textbf{(Count-multiplicativity per node.)} Fix a
node of the tree with order-$r$ residue field $\Fr$ ($|\Fr|=Q_r$), a chosen
$\phi_r$-Newton-polygon side, and a residual-pattern decoration. As the on-side
$\phi_r$-adic coefficients of $f$ range over their residue cells, the order-$r$
residual coefficient map is a \emph{submersion} onto $\mathbb A^{d_r}(\Fr)$ (the
residual $R_r$ is equidistributed over monic degree-$d_r$ polynomials over $\Fr$),
and the count of monic $f$ realizing this node (at precision $N$) \emph{factors} as
\begin{equation}\label{eq:nodefactor}
   \#\{\text{node realized}\}
   \;=\;
   \underbrace{N_{\mathrm{res}}(Q_r)}_{\substack{\text{residual-pattern}\\\text{count over }\Fr}}
   \cdot
   \underbrace{(\text{Newton-box count})}_{\text{a }q\text{-power}}
   \cdot
   \prod_{\text{children}}\#\{\text{child node realized}\}.
\end{equation}
\item\label{ax:dichotomy} \textbf{(Leaf/descend dichotomy and strict descent.)} A
multiplicity-$1$ residual factor is a leaf (no children); a multiplicity-$\mu\ge2$
residual factor of residue degree $f$ descends to a child node which is a cluster
of size $\mu$ over $K_{\delta f}$, with $\mu\le d_r$ \emph{strictly less} than the
parent size off the unique non-progressing self-loop cell (slope denominator
$e=1$, residual $(y-c)^{\,\text{size}}$). The descent terminates: every monic
separable $f$ has a finite OM tree with at most $n$ leaves.
\end{enumerate}
\end{axiom}

\begin{remark}[GMN provenance; what is and is not imported]
Axiom~\ref{ax:montes} is the count/measure shadow of the GMN ``Theorem of the
Polygon'' (Thm.\ 3.1) and ``Theorem of the Residual Polynomial'' (Thm.\ 3.7),
the leaf dichotomy (Cor.\ 3.8), the descent trigger (Lemma 3.11(3)), and finite
termination (Thm.\ 4.18 / Cor.\ 4.19). We import the \emph{combinatorial/counting}
content only: the tree exists and partitions, and the node count factors as
\eqref{eq:nodefactor}. We do \emph{not} import any rationality, uniformity, or
$p$-independence statement --- those are \emph{derived} below. The axiom is
$p$-independent: GMN holds over an arbitrary complete discretely valued field, and
the wild/inseparable case is the \emph{generic} case of the machinery, not an
exception. The single genuine analytic content folded into \ref{ax:multiplicativity}
is the equidistribution/submersion claim (``residual coefficients are independent
leading $\Fr$-digits''); this is GMN's order-$r$ coordinate description and is the
order-$r$ analogue of the order-$1$ statement, which is elementary (see
Remark~\ref{rem:wall}).
\end{remark}

\section{The decomposition theorem (proved)}\label{sec:decomp}

\begin{theorem}[Decomposition into a finite sum over shapes]\label{thm:decomp}
Fix $n$ and a type $\sigma$. The combinatorial OM tree-shapes $T$ with
$\type(T)=\sigma$ form a \emph{finite} set $\mathcal T_{n,\sigma}$, and
\begin{equation}\label{eq:decomp}
  \rhocount(n,\sigma;q)\;=\;\sum_{T\in\mathcal T_{n,\sigma}} C_T(q),
  \qquad
  C_T(q):=\lim_{N\to\infty} q^{-nN}\,\#\,S_T^{(N)} = \mu(\text{stratum of }T).
\end{equation}
\end{theorem}

\begin{proof}
\emph{Finiteness of shapes.} By Axiom~\ref{ax:montes}\ref{ax:dichotomy} every
separable $f$ has $\le n$ leaves and finite depth. All combinatorial data of a
shape are bounded by $n$: the residue degrees satisfy $\prod_i f_i\mid n$; each
side has slope $h/e$ with $e\le n$; each residual has degree $d_r\le n$ so its
factorization pattern is a partition of an integer $\le n$; the number of branch
levels is $\le n-1$ and the number of leaves is $\le n$. A bounded-depth tree with
bounded branching and bounded per-node decorations has finitely many shapes; the
menu $\mathcal T_{n,\sigma}$ is finite (and, as a combinatorial set, the same for
every $p$). [This is the finite, $p$-independent shape menu of L5fix.]

\emph{Partition.} By Axiom~\ref{ax:montes}\ref{ax:partition}, for each $N$ the
strata $\{S_T^{(N)}\}_{T}$ partition the monic degree-$n$ polynomials in
$(\OK/p^N)[x]$, and the type-$\sigma$ locus is the disjoint union of those
$S_T^{(N)}$ with $\type(T)=\sigma$.

\emph{Counting-additivity and the limit.} For each $N$, summing the partition,
\[
  q^{-nN}\#\{\type=\sigma\text{ decided at }N\}
  = \sum_{T:\ \type(T)=\sigma} q^{-nN}\,\#\,S_T^{(N)}
\]
(a \emph{finite} sum, by finiteness of $\mathcal T_{n,\sigma}$; the undecided
remainder $U_N$ is the not-yet-resolved part, $\to0$ by Lemma~\ref{lem:countHaar}).
Each summand $q^{-nN}\#S_T^{(N)}$ converges as $N\to\infty$ by the same bracketing
as in Lemma~\ref{lem:countHaar} applied to the single shape $T$ (the stratum is
decided off a sub-discriminant boundary of measure $\to0$). A finite sum of
convergent sequences converges to the sum of the limits, giving
\eqref{eq:decomp}.
\end{proof}

\begin{remark}
Theorem~\ref{thm:decomp} is the statement that is \emph{assumed} in the current
Lean as the \texttt{DensityFoundation.decomposition} field. Here it is a theorem
from Axiom~\ref{ax:montes} (partition $+$ finite menu) and elementary
counting-additivity. \textbf{What is not yet fully internal:} the \emph{partition}
clause \ref{ax:partition} is part of the imported axiom; in the Lean
\texttt{PadicMeasure} layer the analogous content is the \texttt{omCells}
\emph{exhaustiveness} (that the cell strata reconstruct the residual), which is
also still axiomatic. So Theorem~\ref{thm:decomp} reduces the decomposition to the
single partition statement of Axiom~\ref{ax:montes}; it does not eliminate it.
\end{remark}

\section{The geometric collapse (the crux, proved)}\label{sec:collapse}

We now prove that a fixed \emph{combinatorial} shape $T$, whose \emph{literal} OM
tree is \emph{infinite} (one node has unbounded refinement depth --- the self-loop),
nonetheless has rational stratum-measure $C_T(q)$, via a geometric series.

\subsection{The self-loop and the rescale exponent}

Within a fixed combinatorial shape $T$, the only source of infinite depth is the
\emph{self-loop}: a node whose order-$r$ Newton polygon has slope denominator
$e=1$ and residual $R_r=(y-c)^{\,s}$ (all $s$ roots one notch deeper), so the
``shape'' does not progress --- the deeper refinement re-presents a cluster of the
\emph{same size} $s$. The point of this section is that, although the literal tree
is infinite, the conditional cluster volume $V_s$ is rational; \textbf{but the
collapse is a single geometric series only for $s=2$.} For $s\ge3$ the
``go-one-level-deeper'' region is not a fresh copy of the same node: it splits into
\emph{several} smaller clusters of sizes $b_1,\dots,b_r$ with $\sum b_i=s$,
$r\ge2$, each of which is itself an infinite self-loop of a \emph{strictly smaller}
size. The collapse is therefore a well-founded \emph{nested system} of geometric
series (Theorem~\ref{thm:nested}), with the single-block geometric series as its
$s=2$ instance (Lemma~\ref{lem:s2}).

\begin{lemma}[Per-level self-similar rescale exponent]\label{lem:rescale}
Consider the completely-split size-$s$ cluster, conditioned to total mass $1$, with
volume $V_s$ equal to the self-similar Vandermonde integral
\begin{equation}\label{eq:Idef}
  V_s \;=\; I_s\;:=\;\int_{\Odelta^s}\;\prod_{i<j}\bigl|s_i-s_j\bigr|\;
  \prod_i d\mu(s_i)\qquad(\text{normalized }\mu(\Odelta)=1).
\end{equation}
Translating to the cluster ($s_i\equiv c$) and rescaling $s_i=c+\pi t_i$ once
multiplies the integrand $\prod_{i<j}|s_i-s_j|$ by $Q^{-\binom s2}$, the box
$\{s_i\equiv c\}$ by $Q^{-s}$, and offers $Q$ choices of $c$; the \emph{single-block}
(all-roots-stay-together) part of one refinement level therefore rescales the
conditional measure by
\begin{equation}\label{eq:rfactor}
  r(s)\;=\;Q\cdot Q^{-s}\cdot Q^{-\binom s2}\;=\;Q^{-w(s)},\qquad
  w(s)\;=\;s+\binom s2-1\;=\;\frac{s(s+1)}{2}-1
  \;\;(=2,5,9,14,\dots),
\end{equation}
where $Q=q^\delta$ is the residue cardinality at the node.
\end{lemma}

\begin{proof}
On the level of roots the substitution $s_i=c+\pi t_i$ ($\pi=p$, $\vp(\pi)=1$) is a
measure-preserving bijection of $\Odelta^s$ scaled by $\pi$, so
$d\mu(s_i)=Q^{-1}d\mu(t_i)$ on the box (the box has measure $Q^{-s}$ once the common
residue $c$ is fixed), and $|s_i-s_j|=Q^{-1}|t_i-t_j|$ for each of the $\binom s2$
pairs, i.e.\ $\prod_{i<j}|s_i-s_j|=Q^{-\binom s2}\prod_{i<j}|t_i-t_j|$. Summing over
the $Q$ choices of $c$ gives the displayed product $Q^{1-s-\binom s2}=Q^{-w(s)}$.
The Vieta/Jacobian identity $|\det d\Phi|=|\mathrm{Vandermonde}|$ (so the integrand
$|\disc|^{1/2}=|\mathrm{Vandermonde}|$ pushes forward correctly to coefficient
space) holds in \emph{every} characteristic, because $\det d\Phi$ is an
identically-nonzero $\mathbb Z$-polynomial; this is checked exactly for $n\le5$ in
\texttt{r1\_verify.py}. The exponent $w(s)=s(s+1)/2-1=2,5,9,14,20$ is verified
exactly in \texttt{decomp\_blueprint\_verify.py}.
\end{proof}

\begin{remark}[Scope of $r(s)$]
Equation~\eqref{eq:rfactor} is the rescale of the \emph{single-block} part of one
refinement level only. It is \emph{not} the per-level increment ratio of $V_s$:
that ratio is governed by the \emph{slowest} mode of the nested system below, which
is $Q^{-w(2)}=Q^{-2}$ for every $s\ge2$ (a nested size-$2$ sub-cluster). The two
coincide only at $s=2$, where the single block is the whole loop. This distinction
is the substance of fix~[F1]; see Remark~\ref{rem:slowmode}.
\end{remark}

\subsection{The size-$2$ self-loop: a single geometric series}

\begin{lemma}[$s=2$ collapse, exactly verified]\label{lem:s2}
The size-$2$ completely-split cluster volume satisfies the \emph{single}
self-referential identity
\begin{equation}\label{eq:s2self}
  I_2 \;=\; a + r\,I_2,\qquad a=\frac{q-1}{q},\quad r=q^{-w(2)}=q^{-2},
\end{equation}
whose finite truncations $V_N=a(1-r^{m(N)+1})/(1-r)$ converge (Lemma~\ref{lem:countHaar})
to
\begin{equation}\label{eq:s2limit}
  \boxed{\,I_2=\frac{a}{1-r}=\frac{(q-1)/q}{1-q^{-2}}=\frac{q}{q+1}\,.}
\end{equation}
\end{lemma}

\begin{proof}
Stratify $\Odelta^2$ by the mod-$p$ residue pattern of $(s_1,s_2)$.
\emph{Distinct residues} ($s_1\not\equiv s_2$): the difference is a unit,
$|s_1-s_2|=1$; this contributes the fixed mass $a=\mu\{s_1\not\equiv s_2\}\cdot1
=1-q^{-1}=(q-1)/q$ (a leaf: the cluster has resolved into two distinct roots). This
$a$ is \emph{precision-independent}: it is decided already at $N=1$ and is exactly
$\tfrac12,\tfrac23,\tfrac45$ at $p=2,3,5$ (verified). \emph{Equal residues}
($s_1\equiv s_2\equiv c$): the unique single block; by Lemma~\ref{lem:rescale}
($s=2$) this reproduces $I_2$ scaled by $r=q^{-w(2)}=q^{-2}$. Because for $s=2$
``stay together'' is the \emph{only} non-leaf option, \eqref{eq:s2self} is a genuine
single geometric series, and Lemma~\ref{lem:collapse} below applies verbatim with a
\emph{constant} $a$. Evaluating the limit gives \eqref{eq:s2limit}.
\end{proof}

\begin{lemma}[Single-geometric collapse (the $s=2$ engine)]\label{lem:collapse}
Suppose a node-measure $V$ satisfies $V=a+rV$ with $a$ \emph{constant in $N$} and
$0<r<1$. Let $V_N$ be the contribution of self-loop chains of length $\le m(N)$
(with $m(N)\to\infty$ as $N\to\infty$) together with the resolve-here cells. Then
$V_N=a(1-r^{m(N)+1})/(1-r)$ is a finite geometric sum, and
\begin{equation}\label{eq:limit}
  V \;=\; \lim_{N\to\infty} V_N \;=\; \frac{a}{1-r},
\end{equation}
a rational function of $q$ with denominator $1-r$.
\end{lemma}

\begin{proof}
Unwinding $V=a+rV$ exactly $m+1$ times gives
$V=a(1+r+\dots+r^{m})+r^{m+1}V$. At precision $N$ the deepest self-loop chain the
digits below $p^N$ certify has length $m(N)$ (one self-loop level consumes a bounded
number of $p$-adic digits, so $m(N)\to\infty$), and the tail $r^{m+1}V$ is the
undecided remainder $U_N$ of the bracket in Lemma~\ref{lem:countHaar} restricted to
this node. Thus $V_N=a\sum_{j=0}^{m(N)}r^j=a(1-r^{m(N)+1})/(1-r)$, and as $0<r<1$
the bracket $V_N\le V\le V_N+r^{m(N)+1}V$ forces $V=a/(1-r)$.
\end{proof}

\begin{remark}[Why $s\ge3$ is not this lemma]\label{rem:slowmode}
For $s\ge3$ the hypothesis ``$a$ constant in $N$'' \emph{fails}: the lumped
resolve-here mass acquires contributions from nested sub-clusters that are
themselves only resolved at greater depth, so $a_N$ strictly increases with $N$.
Exact enumeration over $(\mathbb Z/p^N)^3$ makes this unambiguous: at $p=2$ the
graded lower bounds are
$L_N=0,\tfrac3{16},\tfrac{123}{512},\tfrac{4155}{16384},\dots$ (strictly
increasing), and the depth-increment ratio
$D_N/D_{N-1}=\tfrac9{32},\tfrac{73}{288},\tfrac{585}{2336},\tfrac{4681}{18720},
\tfrac{37449}{149792}\to\tfrac14=q^{-2}$, \emph{not} $q^{-w(3)}=q^{-5}=\tfrac1{32}$
(and at $p=3$, $\to\tfrac19=q^{-2}$, not $q^{-5}$). The slowest geometric mode is
the nested size-$2$ self-loop. So the correct object for $s\ge3$ is the
\emph{nested} recursion of \S\ref{sec:nested}, of which $I_2$ is the base case.
\end{remark}

\subsection{The nested multi-block recursion (the crux for general $s$)}\label{sec:nested}

\begin{theorem}[Self-loop collapse as a well-founded nested system]\label{thm:nested}
Let $I_s$ be the conditional size-$s$ completely-split cluster volume
\eqref{eq:Idef} with $Q=q$ ($\delta=1$; the general node replaces $q$ by $Q$). Then,
writing $w(s)=s(s+1)/2-1$ and $\binom{0}{2}=\binom12=0$,
\begin{equation}\label{eq:nested}
  I_s\,\bigl(1-q^{-w(s)}\bigr)
  \;=\;
  \sum_{\substack{P\,\vdash\,\{1,\dots,s\}\\ \#P=r\,\ge\,2}}
    \Bigl[\prod_{j=0}^{r-1}(q-j)\Bigr]
    \;\prod_{B\in P} q^{-\bigl(|B|+\binom{|B|}2\bigr)}\,I_{|B|},
  \qquad I_0=I_1=1,
\end{equation}
the sum running over set-partitions $P$ of $\{1,\dots,s\}$ into $r\ge2$ nonempty
blocks. Equivalently, the single-block ($r=1$, ``all stay together'') term has been
moved to the left, producing the self-loop pivot $1-q^{-w(s)}$; every block on the
right has size $|B|<s$, so \eqref{eq:nested} is a \emph{well-founded} recursion
(strict descent), terminating in the base $I_0=I_1=1$.
\end{theorem}

\begin{proof}
Stratify $\Odelta^s$ by the mod-$p$ residue \emph{partition} of $(s_1,\dots,s_s)$:
which roots share a residue. A partition into $r$ residue-classes of sizes
$|B_1|,\dots,|B_r|$ contributes as follows.
\emph{(Choice of residues.)} The $r$ distinct residues $c_1,\dots,c_r\in\Fq$ can be
chosen (as an ordered-then-deduplicated assignment to the labelled blocks) in
$\prod_{j=0}^{r-1}(q-j)=q(q-1)\cdots(q-r+1)$ ways; this is the falling factorial
$N_{\mathrm{res}}$.
\emph{(Cross-block Vandermonde.)} For $i,j$ in different residue classes
$|s_i-s_j|=1$ (unit difference), so cross-block factors contribute $1$.
\emph{(Within-block self-similarity.)} For each block $B$ of size $b=|B|$, the
within-block roots all share residue $c$; rescale $s_i=c+\pi t_i$. By
Lemma~\ref{lem:rescale}'s computation \emph{without} the $Q$ residue-choice (already
counted above), each block contributes the box measure $q^{-b}$ times the
within-block Vandermonde rescale $q^{-\binom b2}$ times a \emph{fresh} size-$b$
cluster volume $I_b$, i.e.\ a factor $q^{-(b+\binom b2)}I_b$. The blocks are
independent (cross-block differences are units), so the within-stratum measure is
the product over blocks.
\emph{(Single block vs.\ $\ge2$ blocks.)} The $r=1$ term is exactly the
all-stay-together single block: its factor is
$q\cdot q^{-(s+\binom s2)}I_s=q^{-w(s)}I_s$, the self-loop. Moving it to the left
yields the pivot $1-q^{-w(s)}$ on $I_s$. The remaining strata ($r\ge2$) are the
right side of \eqref{eq:nested}. Each such block has $|B|<s$ (since $r\ge2$), so the
$I_{|B|}$ are strictly smaller; the recursion is well-founded with base
$I_0=I_1=1$ ($I_1$: a single root, volume $1$; $I_0$: empty product). Convergence of
the finite-precision truncations to the displayed rational solution follows from
Lemma~\ref{lem:countHaar} applied to each stratum (each is decided off a
sub-discriminant boundary of measure $\to0$), and the system has finitely many
strata at each level. $\square$
\end{proof}

\begin{corollary}[Closed forms, exactly verified]\label{cor:closed}
Solving \eqref{eq:nested} bottom-up gives
\[
  I_2=\frac{q}{q+1},\qquad
  I_3=\frac{q^3(q^2-q+1)}{(q+1)\,\Phi_5},\qquad
  I_4=\frac{q^6\,(q^8-2q^7+q^6+2q^5-q^4+2q^3+q^2-2q+1)}
          {(q+1)^2(q^2+q+1)(q^6+q^3+1)\,\Phi_5},
\]
$\Phi_5=q^4+q^3+q^2+q+1$. Each is rational with denominator a product of cyclotomic
factors $q^{w(s')}-1$, $s'\le s$ (e.g.\ $I_3$ carries the $w(2)$-pole $q+1$ and the
$w(3)$-pole $\Phi_5\mid q^5-1$).
\end{corollary}

\paragraph{Exact verification (no floats).}
The script \texttt{notes/decomp\_blueprint\_verify.py} runs three exact
(rational-arithmetic) blocks at $p\in\{2,3,5,7\}$ (wild $p=2$ included):
\begin{itemize}[nosep]
\item \textbf{(A) $s=2$ single geometric series} (Lemma~\ref{lem:s2}): the
depth-graded enumeration over $(\mathbb Z/p^N)^2$ has constant $a=(q-1)/q$, per-level
ratio \emph{exactly} $r=q^{-2}=p^{-w(2)}$ ($\tfrac14,\tfrac19,\tfrac1{25}$), and tail
$a/(1-q^{-2})=q/(q+1)$ ($\tfrac23,\tfrac34,\tfrac56$), inside the rigorous bracket
$[L_N,L_N+U_N]$.
\item \textbf{(B) $s=3$ single-geometric model fails} (Remark~\ref{rem:slowmode}):
$a_N$ strictly increases and $D_N/D_{N-1}\to q^{-2}$ (not $q^{-w(3)}$). This is shown
deliberately, to motivate the nested recursion.
\item \textbf{(C) nested recursion} (Theorem~\ref{thm:nested}): the bottom-up
solution reproduces $I_2,I_3,I_4$ of Corollary~\ref{cor:closed} symbolically, and
the closed $I_3$ lies inside the \emph{rigorous} exact bracket from direct
$(\mathbb Z/p^N)^3$ enumeration: $I_3=8/31\in[0.25800,0.28131]$ at $p=2$, and
$189/484\in[0.39026,0.42699]$ at $p=3$.
\end{itemize}
The symbolic backbone of (C) is \texttt{notes/r1\_verify.py} (the \texttt{I\_recursion},
which is bottom-up and non-circular --- base $I_0=I_1=1$, no closed-form reference).

\begin{remark}[Relation to the Lean lemma]
The \emph{algebraic} fact $\sum_{j\ge0}r^j=1/(1-r)$ for $|r|<1$ and the positivity of
the pivot $1-q^{-w}$ for $q\ge2$, $w\ge1$ is exactly
\texttt{L5fix.selfLoop\_geometric} (proved, sorry-free); it is the engine of
Lemma~\ref{lem:collapse} and the pivot in Theorem~\ref{thm:nested}. The blueprint's
new content is twofold: (i)~the \emph{measure/counting} upgrade --- the partial sums
$V_N$ are honest finite counts and the limit uses the bracket of
Lemma~\ref{lem:countHaar}, not merely the algebraic identity --- and (ii)~the correct
\emph{nesting}: the single pivot is applied to a finite $\ge2$-block sum of
strictly-smaller child volumes, matching the well-founded
\texttt{OMInduction.clusterVol\_isRational} (not a single self-referential equation
for $s\ge3$).
\end{remark}

\section{Rationality of each $C_T$ and of $\rho$ (proved)}\label{sec:rational}

\begin{proposition}[The node ingredients are rational]\label{prop:ingredients}
At any node, the three factors of \eqref{eq:nodefactor} are uniform rational
functions of $q$ with $p$-independent data:
\begin{enumerate}[label=\textup{(\alph*)},nosep]
\item the \textbf{Newton-box count} is a $q$-power $q^{-A}$, $A$ pure lattice data
of the side, independent of $p$ (a wild slope $h/e$ with $p\mid e$ gives the same
lattice polygon, hence the same exponent, as a tame slope of the same reduced
$(h,e)$);
\item the \textbf{residual-pattern count} $N_{\mathrm{res}}(Q_r)$ is a single
polynomial in $Q_r=q^{\delta f_1\cdots f_{r-1}}$ with $p$-independent integer
coefficients (over the finite field $\Fr$: the count of monic degree-$d$
polynomials of any fixed factor-multiplicity pattern is a Gauss/M\"obius polynomial
in $Q_r$; the non-squarefree/descent locus has count $Q_r^{d-1}$). Finite fields
are perfect, so there are \emph{no inseparable-irreducible exceptions in any
characteristic}: ``needs descent $=$ not squarefree $=$ $\{\disc R_r=0\}$'';
\item the \textbf{self-loop pivot} $1/(1-Q^{-w(s)})$ is rational with denominator
$Q^{w(s)}-1$, pole-free at every prime power $q\ge2$ (Lemma~\ref{lem:collapse},
Theorem~\ref{thm:nested}); for $s\ge3$ it multiplies a finite sum over $\ge2$-block
partitions of strictly-smaller child volumes, not a single self-loop.
\end{enumerate}
\end{proposition}

\begin{proof}
(a) The Newton-box count is the number of residue classes (mod $p^N$) of on-side
$\phi_r$-adic coefficients satisfying the valuation conditions $\vp(\cdot)\ge
\lceil\text{height}\rceil$ (a ball) or $\vp(\cdot)=\text{height}$ (a shell);
by \eqref{eq:index} restated as counting (Remark~\ref{rem:wall}) this is
$(1-q^{-1})^V q^{-A}$ with $V,A$ the lattice point-counts of the polygon side --
pure integer lattice data depending only on the reduced $(h,e)$, hence
$p$-independent. (b) The squarefree generating function over $\mathbb F_{Q_r}$ is
$\zeta(s)/\zeta(2s)=(1-Q_rt^2)/(1-Q_rt)$ in $t=Q_r^{-s}$, giving $Q_r^d-Q_r^{d-1}$
squarefree and $Q_r^{d-1}$ non-squarefree monic of degree $d\ge2$; fine patterns
are products of binomials in the M\"obius irreducible-count
$I_d(Q_r)=\frac1d\sum_{e\mid d}\mu(e)Q_r^{d/e}$ --- all polynomials in $Q_r$ with
$p$-independent integer coefficients; perfectness of $\Fr$ removes the only
possible characteristic-dependent exception. (c) Is Lemma~\ref{lem:collapse}.
\end{proof}

\begin{theorem}[Rationality]\label{thm:rational}
For every shape $T$ the stratum-measure $C_T(q)$ is a uniform rational function of
$q$ (one element of $\mathbb Q(t)$, valid at every prime power $q$ and every
residue characteristic $p$, wild $p\le n$ included), and hence
$\rhocount(n,\sigma;q)=\sum_{T}C_T(q)$ is a uniform rational function of $q$.
\end{theorem}

\begin{proof}
By induction on the node size $s$ (well-founded by
Axiom~\ref{ax:montes}\ref{ax:dichotomy} and Theorem~\ref{thm:nested}: off the
self-loop every block is strictly smaller, and the single self-loop block is
resummed via the pivot $1-Q^{-w(s)}$, not recursed). \emph{Base:} a leaf is a
multiplicity-$1$ residual factor; its contribution is a Newton-box $q$-power times
a separable-residual count, both rational by Proposition~\ref{prop:ingredients}(a,b),
no recursion. \emph{Step:} by the node factorization \eqref{eq:nodefactor} together
with the self-loop pivot of Theorem~\ref{thm:nested},
\[
  C_T(q)
  = \frac{1}{1-Q^{-w(s)}}
    \sum_{\substack{\text{resolve-here cells}\\(\ge2\text{ blocks / branch / descend})}}
    \Bigl[(1-q^{-1})^{V}q^{-A}\,N_{\mathrm{res}}(Q_r)
      \prod_{\text{children (size}<s)} C_{T_{\text{child}}}(q)\Bigr],
\]
a \emph{finite} (Theorem~\ref{thm:decomp}) $\mathbb Q(t)$-combination of:
$q$-powers (Prop.~\ref{prop:ingredients}(a)), residual-pattern polynomials in $Q_r$
(Prop.~\ref{prop:ingredients}(b)), strictly-smaller child measures (rational by the
inductive hypothesis), and the pole-free self-loop pivot
(Prop.~\ref{prop:ingredients}(c)). The displayed recursion is exactly
\eqref{eq:nested} in the completely-split model and its node-level generalization
otherwise. $\mathbb Q(t)$ is closed under finite sums,
products, and division by a factor that is pole-free at every prime power; hence
$C_T\in\mathbb Q(t)$. Summing the finite menu (Theorem~\ref{thm:decomp}) keeps
$\rhocount(n,\sigma;\cdot)\in\mathbb Q(t)$.
\end{proof}

\begin{remark}[The wild-prime crux, in one line]
The only place a characteristic could intrude is the residual-pattern count
$N_{\mathrm{res}}$. There the descent-trigger locus
$\{R_r\ \text{not squarefree}\}$ may be a \emph{geometrically different
subvariety} at a wild prime (e.g.\ at $p=2$ the double-root locus is the line
$\{\bar b=0\}$, at odd $p$ the parabola $\{\bar b^2=4\bar a\}$ --- the very data
that broke the referee's Prop.\ 5.3, a tangent-space statement) --- \emph{but it
has the same $\Fr$-point count} because the merge map (Frobenius and its powers)
is bijective on $\Fr$-points even when its differential vanishes. Counts/volumes
live on points; smoothness lives on tangent spaces. This is why the collapse and
the decomposition are characteristic-uniform.
\end{remark}

\section{What still needs more than Montes $+$ counting (honest)}\label{sec:honest}

We separate (A) the imported axiom, (B) what is proved from it, and (C) what
genuinely needs an extra hypothesis or is verified only in restricted scope.

\paragraph{(A) Imported (axiom).} Axiom~\ref{ax:montes}: the tree exists and
\emph{partitions} the polynomial space \ref{ax:partition}; the node count
\emph{factors} as \eqref{eq:nodefactor}, including the \emph{submersion /
equidistribution} of residual coefficients onto $\mathbb A^{d_r}(\Fr)$
\ref{ax:multiplicativity}; the leaf/descend dichotomy with strict descent and
finite termination \ref{ax:dichotomy}. This is GMN ($p$-independent).

\paragraph{(B) Proved from (A) $+$ elementary counting.}
\begin{itemize}[nosep]
\item Finiteness of the shape menu $\mathcal T_{n,\sigma}$ (Thm.~\ref{thm:decomp});
\item the decomposition $\rho=\sum_T C_T$ via counting-additivity over the
partition and the $N\to\infty$ limit (Thm.~\ref{thm:decomp},
Lem.~\ref{lem:countHaar});
\item the per-level rescale exponent $w(s)=s(s+1)/2-1$ (Lem.~\ref{lem:rescale});
the $s=2$ single-geometric collapse
$V_N=a(1-r^{m+1})/(1-r)\to a/(1-r)=q/(q+1)$ with constant $a=(q-1)/q$,
$r=q^{-2}$ (Lem.~\ref{lem:s2},~\ref{lem:collapse}); and \textbf{for $s\ge3$ the
well-founded nested multi-block recursion} \eqref{eq:nested}
$I_s(1-q^{-w(s)})=\sum_{\ge2\text{ blocks}}\prod_j(q-j)\prod_B q^{-(|B|+\binom{|B|}2)}I_{|B|}$,
giving $I_3,I_4$ of Cor.~\ref{cor:closed} (Thm.~\ref{thm:nested});
\item rationality of each $C_T$ and of $\rho$ (Prop.~\ref{prop:ingredients},
Thm.~\ref{thm:rational}).
\end{itemize}

\paragraph{(C) Where Montes $+$ counting alone may not suffice --- flagged honestly.}
\begin{enumerate}[label=(C\arabic*),nosep]
\item\textbf{The self-loop rescale exponent $w(s)=s(s+1)/2-1$ and factor $r$
(Lemma~\ref{lem:rescale}) is \emph{not} part of Axiom~\ref{ax:montes}.} It is the
R1 self-similar Vandermonde/Jacobian computation (a change-of-variables on the
completely-split root locus). Axiom~\ref{ax:montes} gives only the
\emph{existence} of the self-loop cell and that it does not progress; the precise
ratio $Q^{-w(s)}$ is an \emph{additional} elementary computation (the $p$-adic
change of variables $|\!\det d\Phi|=|\mathrm{Vandermonde}|$, valid in every
characteristic since $\det d\Phi$ is an identically-nonzero $\mathbb Z$-polynomial).
We have verified $w(s)$ exactly for $s=2$ by direct enumeration ($p=2,3,5$) and
for $s=2,3,4$ symbolically; it should be treated as a \emph{separate} proved
input, not folded into the Montes axiom. For an \emph{irreducible/ramified}
(non-completely-split) cluster the literal $\OK^s$ Vandermonde pushforward is a
measure-zero conjugate locus and does \emph{not} compute the volume directly; its
measure is delivered by the Newton-box count (a) plus the recursion, with R1
contributing only the lattice quantity. This scoping is load-bearing.
\item\textbf{The equidistribution/submersion clause inside
Axiom~\ref{ax:montes}\ref{ax:multiplicativity}} (residual coefficients are
independent uniform $\Fr$-digits) is the one genuinely analytic ingredient bundled
into the axiom. At order $1$ it is elementary (it is the box-count of (a)). At
higher order it is GMN's order-$r$ coordinate description; if one wished to
\emph{avoid} importing it, one would re-derive the higher key-polynomial
coordinates from MacLane--Okutsu--Montes valuation theory --- standard but
lengthy, and $p$-independent. As stated it is imported, not proved.
\item\textbf{The collapse is a single geometric series only for $s=2$
(structural, fix~[F1]).} For $s\ge3$ the size-$s$ self-loop is \emph{not} a single
self-referential $V=a+rV$ with constant $a$ and a single ratio $r=q^{-w(s)}$: that
model is \emph{false}, as exact enumeration shows (the lumped $a_N$ strictly
increases with precision; the depth-increment ratio converges to the slowest mode
$q^{-2}$, not $q^{-w(s)}$ --- Remark~\ref{rem:slowmode}). The correct object is the
nested recursion \eqref{eq:nested}, a finite system of interleaved geometric series
with pivots $q^{-w(b)}$, $b\le s$. \emph{Direct} exact enumeration over
$(\mathbb Z/p^N)^s$ confirms the $s=2$ single-geometric collapse exactly and gives
a \emph{rigorous bracket} $[L_N,L_N+U_N]$ containing the recursion's $I_3$
($8/31\in[0.25800,0.28131]$ at $p=2$; $189/484\in[0.39026,0.42699]$ at $p=3$); but
brute force does not reach the asymptotic regime to \emph{pin} $I_3,I_4$ as
rationals, so those rest on the \emph{symbolic} nested recursion (R1), not on direct
enumeration. The empirical OM-tree oracle reaches order $4$ ($n=8$). Genuinely
arbitrary depth rests on the induction (Thm.~\ref{thm:rational}) $+$
Axiom~\ref{ax:montes}, not on enumeration.
\item\textbf{The displayed $\rhocount(3,1^3;q)=(q^2+1)^2/(6\Phi_5)$ is a
\emph{projective} density, not the monic counting density of
Definition~\ref{def:count}.} It uses the symmetry factor $3!=6$ and division by
$|\mathbb P^3|=(q+1)(q^2+1)$ (the $r$-point / distinct-roots normalization), and
lies \emph{outside} the exact monic brackets ($0.1344$ vs.\ monic
$\in[0.169,0.180]$ at $p=2$). The monic type densities of Definition~\ref{def:count}
are different rationals --- e.g.\ the inert quadratic density is exactly $(q-1)/(2q)$
($\tfrac14,\tfrac13,\tfrac25,\tfrac37$ at $p=2,3,5,7$), the ramified quadratic
$1/(q+1)$, the split quadratic $(q^2+1)/(2q(q+1))$ (all exactly enumeration-verified).
The blueprint's theorems (decomposition, collapse, rationality) hold for either
normalization; only the \emph{worked numeric assembly} must be labeled projective.
This blueprint exhibits its Definition~\ref{def:count} monic $\rhocount$ concretely
only at $n=2$ (the verified inert/ram/split forms above); the $n=3$ projective
number is illustrative, not a verified instance of Definition~\ref{def:count}.
\item\textbf{The partition clause \ref{ax:partition} is imported, not derived.}
Theorem~\ref{thm:decomp} reduces the decomposition to this single statement (the
strata cover and are disjoint); in the counting model disjointness is automatic
(unique tree), but \emph{exhaustiveness} (the strata reconstruct the full residual
factorization at every order) is the GMN cell-decomposition content and remains
axiomatic. This matches the Lean status: \texttt{omCells} exhaustiveness is still a
structural axiom.
\end{enumerate}

\section{Formalization map}\label{sec:formalization}

We record what each piece corresponds to in the current Lean development
(\texttt{forum-sigma/lean\_urat/LeanUrat/}) and what this blueprint would let us
\emph{derive} versus keep as axiom.

\paragraph{Already in Lean (proved, sorry-free).}
\begin{itemize}[nosep]
\item \texttt{RatFn.IsRationalFn} and its closure lemmas
(\texttt{isRationalFn\_const/qpow/add/mul/listSum/listProd/div}) --- the
$\mathbb Q(t)$-closure used in Theorem~\ref{thm:rational}; in particular
\texttt{isRationalFn\_qpow} represents a residual count $q^{\delta(d-1)}$ by the
\emph{non-constant} numerator $X^{\delta(d-1)}$.
\item \texttt{L4.bb1Value}/\texttt{L4.cellVolume\_eq} (= ``T\_BB1''): the
Newton-box volume $(1-q^{-1})^Vq^{-A}$ is a rational function of $q$
(Prop.~\ref{prop:ingredients}(a)), \emph{derived} from the column-measure axiom.
\item \texttt{PadicMeasure.T\_BB3} / \texttt{L3Squarefree} (= ``T\_BB3''): the
non-squarefree count over a finite field is $|F|^{d-1}$
(Prop.~\ref{prop:ingredients}(b)) --- a real finite-field theorem, not an axiom.
\item \texttt{L5fix.selfLoop\_geometric}: the pivot $1-q^{-w(s)}>0$ for $q\ge2$,
$s\ge2$ (positivity used in Lemma~\ref{lem:collapse}); the \emph{algebraic}
geometric-series fact.
\item \texttt{OMInduction.clusterVol\_isRational}: the abstract well-founded
OM-order induction --- given a recursion of the shape \eqref{eq:nodefactor}$/$pivot
with strictly-smaller children, rational coefficients, and a pole-free pivot, the
solution is a uniform rational function of $q$. This \emph{is}
Theorem~\ref{thm:rational}'s induction, already formalized.
\end{itemize}

\paragraph{Now PROVED in Lean by the box-wise refactor (was ``would derive'').}
The previous draft listed these as \emph{assumed} in Lean; the counting-model and
box-wise refactor have since \emph{discharged} them. \S\ref{sec:ledger} gives the
exact ledger; in summary:
\begin{itemize}[nosep]
\item \texttt{DensityFoundation.decomposition}\,/\,the decomposition
$\rho=\sum_T C_T$ is now the \emph{proved} theorem
\texttt{MontesData.countingDensity\_eq\_sum\_coeff} (core-only \texttt{\#print
axioms}), a genuine $N\to\infty$ limit-uniqueness theorem (Thm.~\ref{thm:decomp}),
not a carried field.
\item The cell/measure recursion \eqref{eq:nodefactor}$/$pivot is now the
\emph{proved} theorem \texttt{MontesData.nodeMultiplicativity} (core-only), derived
from the box-wise primitives below. The opaque \texttt{PadicMeasure.AX\_cellRecursion}
/ \texttt{clusterMeasure} axioms are \emph{off the path entirely} --- they live in
the dead \texttt{PadicMeasure} module and do \emph{not} appear in the
\texttt{\#print axioms} footprint of any rationality/decomposition constant.
\end{itemize}

\paragraph{New pieces the formalization would need.}
\begin{itemize}[nosep]
\item \textbf{The counting model.} Definition~\ref{def:count} and
Lemma~\ref{lem:countHaar} (normalized counts over $\OK/p^N$ converge; the
bracket $L_N\le\mu\le L_N+U_N$ with $U_N\to0$). This is elementary but new ---
it is what removes the abstract $p$-adic Haar wall (mathlib has no $p$-adic Haar
measure).
\item \textbf{The self-loop collapse as a measure/counting limit.} For $s=2$:
$V_N=a(1-r^{m+1})/(1-r)\to a/(1-r)$ on the \emph{finite-count} truncations
(Lemma~\ref{lem:collapse}), upgrading the algebraic \texttt{selfLoop\_geometric} to
the actual limit of stratum counts. For $s\ge3$: the \textbf{well-founded nested
recursion} \eqref{eq:nested} (Theorem~\ref{thm:nested}), in which the single pivot
$1/(1-q^{-w(s)})$ multiplies a finite $\ge2$-block sum of \emph{strictly-smaller}
child volumes $I_{|B|}$ --- this is exactly the input shape of
\texttt{OMInduction.clusterVol\_isRational}, so the existing Lean induction consumes
it directly.
\item \textbf{The per-level rescale exponent} $w(s)=s(s+1)/2-1$ and the single-block
factor $r(s)=q^{-\delta w(s)}$ (Lemma~\ref{lem:rescale}), the R1 Vandermonde/Jacobian
input --- a separately proved elementary computation, \emph{not} part of the Montes
axiom (flag C1).
\end{itemize}

\paragraph{What stays axiom vs.\ becomes theorem.}
\emph{Stays axiom (A):} Montes count-multiplicativity over the OM tree
(Axiom~\ref{ax:montes}) --- partition, node-count factorization with residual
equidistribution, leaf/descend dichotomy, finite termination (GMN,
$p$-independent). \emph{Becomes theorem (B):} the decomposition
$\rho=\sum_T C_T$; the self-loop collapse --- a single geometric series $C_T=a/(1-r)$
for $s=2$ (Lem.~\ref{lem:s2}) and the well-founded nested recursion
\eqref{eq:nested} for $s\ge3$ (Thm.~\ref{thm:nested}); and rationality of every
$C_T$ and of $\rho$ --- all from (A) $+$ elementary counting. The ``sum over the
tree, including the self-loop collapse'' is thereby a \textbf{theorem, not an
assumption}.

\section{What the Lean formalization ASSUMES vs.\ PROVES (exact ledger)}\label{sec:ledger}

This section is the \textbf{authoritative, item-by-item} statement of what the
current Lean development in \texttt{forum-sigma/lean\_urat/LeanUrat/} \emph{assumes}
versus \emph{proves}, written to match the \texttt{\#print axioms} reality exactly
(re-verified 2026-06-21 by \texttt{lake env lean} on each qualified constant). The
distinction is enforced mechanically: ``proved'' means a \texttt{theorem} whose
\texttt{\#print axioms} contains \emph{only} the three Lean-core axioms
$\{\texttt{propext},\ \texttt{Classical.choice},\ \texttt{Quot.sound}\}$ and
\emph{no} \texttt{sorryAx}; ``assumed'' means a declared \texttt{axiom} or a
\emph{hypothesis} (structure field / theorem argument) that the conclusion is stated
modulo. ``Sorry-free modulo axioms'' is \emph{not} the same as ``proved'': we mark
which axioms each headline depends on.

\subsection{The assumed ledger (exhaustive)}\label{sec:assumed}

Beyond the three Lean-core axioms, the \emph{only} things assumed are:

\begin{enumerate}[label=\textup{(A\arabic*)},leftmargin=*]
\item\textbf{The minimal box-Haar / equidistribution axiom} --- the single
irreducible measure residue, carried as \emph{two structure fields} of
\texttt{MontesAxiom.MontesData} (\emph{not} a global \texttt{axiom}; it is
hypothesis-discharged, so it does not surface in \texttt{\#print axioms}):
\begin{itemize}[nosep]
\item \texttt{boxMeasure : CountCell $\to\ \mathbb N\to\mathbb Q$} --- the abstract
per-box ($p$-adic Haar) volume of one cell;
\item \texttt{boxHaarEquidist} : $\forall c\ q'>1,\
\texttt{boxMeasure}\ c\ q' = \texttt{countCellCoeff}\ c\ q'$, i.e.\ the per-box Haar
volume \emph{equals} the proved box-volume $\times$ residual-count value
$\texttt{countCellCoeff}\ c\ q = (q^{\delta})^{d_S-1}\cdot\texttt{bb1Value}\ c.\texttt{polygon}\ q$;
\item \texttt{nodeMeasure\_boxSum} : the normalized stratum limit
$\lim_N \texttt{stratumCount}\ T\ N / q^{nN}$ \emph{exists} and equals the
box-additive sum $\bigl(\sum_{\text{cells }c}\texttt{boxMeasure}\ c\ q\cdot
\prod_{\text{children}}\texttt{clusterCount}\dots\bigr)/\texttt{countPivot}$.
\end{itemize}
\textbf{What is assumed} is exactly: that the abstract per-box Haar volume \emph{is}
the proved closed value (box-Haar normalization $T\_BB1$ $+$ C2 residual
equidistribution $T\_BB3$), and that the node measure \emph{exists} and is the
box-additive assembly. The right-hand sides ($\texttt{bb1Value}$ box volume, the
$(q^{\delta})^{d_S-1}$ residual count, and the box-additivity \emph{arithmetic}) are
all \emph{proved} (see \S\ref{sec:proved}); only the \emph{identification} of the
abstract Haar volume with that proved value is assumed. This is the genuine
\emph{measure wall}: mathlib v4.31.0 has no $p$-adic Haar measure, so the existence
and normalization of \texttt{boxMeasure} cannot be constructed inside Lean. It is
\emph{not} a renamed conclusion and \emph{not} a definitional-unfold cheat: the
asserted equality is to a separately-proved lattice/finite-field value, and the
fields assert \emph{no} rationality, uniformity, or palindromy.

\item\textbf{The remaining \texttt{MontesData}/\texttt{CountingModel} structure
fields}, all carried as \emph{hypotheses} (the GMN count shadow), each the
count-native shadow of a GMN statement, none of them a rationality/conclusion claim:
\begin{itemize}[nosep]
\item \texttt{shapesOf}, \texttt{partition} ($\delta$-A:partition, the unique-tree
disjoint partition $\texttt{decidedCount}\ \sigma\ N = \sum_T \texttt{stratumCount}\ T\ N$);
\item \texttt{treeSize}, \texttt{cells}, \texttt{cells\_descend} (strict descent
--- the count shadow of \texttt{omCells}/\texttt{descend\_size\_lt}, supplying
well-foundedness; carried as fields, \emph{not} the \texttt{PadicMeasure} axioms);
\item \texttt{finiteTermination} (GMN Thm 4.18, $\le n$ leaves);
\item the elementary tail-bound fields \texttt{undecided\_le\_split} ($\le$),
\texttt{sepTail\_tendsto} (separable pool $\to0$), \texttt{discriminantTail\_envelope}
(the $q^{-N}$ root-count envelope) --- elementary counting bounds, \emph{not} the
conclusion $U_N\to0$ (which is \emph{derived}, see (P4)).
\end{itemize}
None of these carries rationality, uniformity, or palindromy. There is no
``\texttt{undecided}$\to0$'' field and no ``$C_T$ is rational'' field: those are
theorems (\S\ref{sec:proved}).

\item\textbf{\texttt{tame\_functionalEquation}} (\texttt{Interface.lean}, a declared
\texttt{axiom}) --- the genuine \S5 H-tame input: an abstract
\texttt{DensityFoundation}'s density admits a palindromic rational representative at
tame primes. It is the \emph{only} non-core declared axiom in the headline footprint,
and it appears \emph{only} on \texttt{goal\_theorem\_montes} (and the witness
instance), \emph{not} on any pure-rationality constant. It asserts nothing about the
count-native $\sum_T C_T$ directly; the link is made through the explicit honest
hypothesis \texttt{hbridge} (below).

\item\textbf{\texttt{hbridge}} (a hypothesis argument of
\texttt{goal\_theorem\_montes}, \emph{not} an axiom): $\forall q'>1,\
F.\texttt{density}\ n\ \sigma\ q' = \sum_{T}\texttt{coeff}\ T\ q'$. The honest
measure-wall identification ``the abstract foundation density coincides with the
count-native OM tree-sum''. The right-hand side is an arbitrary function here; the
equality asserts \emph{nothing} about rationality, palindromy, or any closed form. It
introduces no axiom.
\end{enumerate}

\noindent\textbf{What is NOT assumed (verified absent from every audited footprint).}
The \texttt{PadicMeasure} measure-route axioms ---
\texttt{clusterMeasure}, \texttt{AX\_cellRecursion}, \texttt{AX\_columnMeasure},
\texttt{omCells}, \texttt{descend\_size\_lt} --- are genuine \texttt{axiom}
declarations in the (dead) \texttt{PadicMeasure} module, and are \emph{absent} from
the \texttt{\#print axioms} of \texttt{goal\_theorem\_montes} and of every
rationality/decomposition constant. No \texttt{sorryAx} appears in any audited
footprint (the two flagged \texttt{sorry}s at \texttt{L3.lean:145,168} are off-path
--- the path uses the sorry-free \texttt{L3Squarefree.card\_squarefreeMonicDegree}
and the sorry-free \texttt{L3.card\_monicDegree}; confirmed by the clean footprints).

\subsection{The proved ledger (theorems, core-only \texttt{\#print axioms})}\label{sec:proved}

Each item below is a Lean \texttt{theorem} whose \texttt{\#print axioms} is exactly
$\{\texttt{propext},\ \texttt{Classical.choice},\ \texttt{Quot.sound}\}$ (no
\texttt{sorryAx}, no measure axiom, no \texttt{tame\_functionalEquation}), verified
2026-06-21.

\begin{enumerate}[label=\textup{(P\arabic*)},leftmargin=*]
\item\textbf{Box-volume $T\_BB1$.} \texttt{MontesAxiom.boxVolume\_eq} (re-exposing
\texttt{L4.cellVolume\_eq}): the Newton-box counting/Haar volume of one cell equals
the closed lattice value $\texttt{bb1Value}\ pg\ Q = (1-Q^{-1})^{V(pg)}\,Q^{-A(pg)}$,
given the per-coordinate box/shell factorization (the one lattice-count input of BB1).
A genuine \texttt{Finset.prod} collapse, \emph{not} re-axiomatized.

\item\textbf{Residual-config count $T\_BB3$.} \texttt{MontesAxiom.residualBoxCount}:
over a finite field $F$ with $|F|=Q$, the count of non-squarefree monic degree-$d_S$
polynomials is $Q^{d_S-1}$, proved \emph{inline} from the honest finite-field counts
\texttt{L3.card\_monicDegree} ($\#\text{monic}=Q^{d_S}$, sorry-free) and
\texttt{L3Squarefree.card\_squarefreeMonicDegree}
($\#\text{squarefree}=Q^{d_S}-Q^{d_S-1}$, sorry-free re-proof). The tie
\texttt{residualBoxCount\_eq\_factor} identifies it with the $(q^{\delta})^{d_S-1}$
factor of \texttt{countCellCoeff} when $|F|=q^{\delta}$. A genuine finite-field count,
\emph{not} a free monomial and \emph{not} re-axiomatized.

\item\textbf{Box-additivity.} \texttt{MontesAxiom.clusterCount\_boxSum} (with
\texttt{clusterCount\_rec}): if a per-box measure agrees on each cell with
$\texttt{countCellCoeff}\cdot\prod_{\text{children}}$, then \texttt{clusterCount} is
the finite \texttt{List.sum} of those per-box measures over the cell partition,
divided by the self-loop pivot. Pure \texttt{List.sum} arithmetic
(\texttt{List.map\_congr\_left}); \emph{not} a measure assumption.

\item\textbf{The per-node factorization \eqref{eq:nodefactor} (now a THEOREM).}
\texttt{MontesAxiom.MontesData.nodeMultiplicativity}: $\lim_N\texttt{stratumCount}\ T\ N/q^{nN}
= \texttt{clusterCount}\dots T\ q$. \emph{Derived} from (P1)$\times$(P2)$\times$(P3)
$+$ the minimal axiom (A1): \texttt{nodeMeasure\_boxSum} gives the box-additive limit;
\texttt{boxHaarEquidist} rewrites each per-box Haar measure to the proved
$\texttt{countCellCoeff}$ value; \texttt{clusterCount\_boxSum} recombines into
\texttt{clusterCount}. This was previously one opaque carried field; it is now a
core-only theorem.

\item\textbf{The geometric collapse / OM induction.}
\texttt{OMInduction.clusterVol\_isRational}: the abstract well-founded
strong-induction over \texttt{treeSize} --- given the divided sum-of-products
recursion with strictly-smaller children, rational coefficients, and a pole-free
pivot, the solution is a uniform rational function of $q$ (Thm.~\ref{thm:rational}'s
induction, fully proved). The count-native ingredients
\texttt{countCellCoeff\_isRational} (residual count $X^{\delta(d_S-1)}$ $\times$
\texttt{bb1Value\_isRational}) and \texttt{countPivot\_isRational} /
\texttt{countPivot\_ne} feed it.

\item\textbf{$U_N\to0$ (undecided vanishes).}
\texttt{MontesAxiom.MontesData.undecidedVanishes}: \emph{derived} by a squeeze
(\texttt{tendsto\_of\_tendsto\_of\_tendsto\_of\_le\_of\_le'}) from
\texttt{undecided\_le\_split} ($\le$), \texttt{sepTail\_tendsto} ($\to0$), and the
$q^{-N}$ \texttt{discriminantTail\_envelope} ($\to0$ via
\texttt{tendsto\_pow\_atTop\_nhds\_zero\_of\_lt\_one}). \emph{Not} assumed.

\item\textbf{The decomposition $\rho=\sum_T C_T$.}
\texttt{MontesAxiom.MontesData.countingDensity\_eq\_sum\_coeff} (Thm.~\ref{thm:decomp}):
$M.\texttt{countingDensity}\ \sigma = \sum_{T\in\texttt{shapesOf}\ \sigma}\texttt{coeff}\ T\ q$,
proved by \texttt{tendsto\_nhds\_unique} between the decided limit
(\texttt{density\_isLimit}, the genuine $N\to\infty$ decided density, \emph{not} the
box count) and the finite-sum limit from \texttt{partition} $+$
\texttt{nodeMultiplicativity}. A genuine limit-interchange theorem, \emph{not} a
\texttt{rfl} tautology; \texttt{countingDensity} is the carried decided limit
(\texttt{densityVal}, pinned by \texttt{decided\_tendsto}), \emph{not} defined as
$\sum_T C_T$.

\item\textbf{Rationality of every $C_T$ and of $\rho$.}
\texttt{MontesAxiom.MontesData.coeff\_isRational} (via the generic engine (P5)) and
\texttt{countingDensity\_isRational} (finite-sum closure). The faithful
\texttt{RatFn.IsRationalFn} (numerators genuinely $q$-vary: residual count
$X^{\delta(d_S-1)}$, box value $(X-1)^V/X^{V+A}$).

\item\textbf{Palindromy transfer.} \texttt{L7.tame\_to\_all\_primes} and
\texttt{L7.isPalindromic\_of\_agree} (the $\mathbb Q(t)$ identity theorem,
\texttt{Polynomial.eq\_zero\_of\_infinite\_isRoot}, degree-robust
\texttt{Interface.IsPalindromic} $R(1/x)=R(x)$) transfer the tame functional equation
onto the \emph{same} count-native $\texttt{num}/\texttt{den}$. Both lemmas core-only.
\end{enumerate}

\subsection{The headline footprints (verbatim \texttt{\#print axioms})}\label{sec:footprints}

\begin{itemize}[nosep]
\item \texttt{Goal.goal\_theorem\_montes} $\to$
$\{\texttt{propext},\ \texttt{Classical.choice},\ \texttt{tame\_functionalEquation},\
\texttt{Quot.sound}\}$ --- Lean core $+$ the one cited H-tame axiom. \emph{No}
\texttt{clusterMeasure}, \texttt{AX\_cellRecursion}, \texttt{omCells},
\texttt{descend\_size\_lt}, or \texttt{sorryAx}.
\item \texttt{countingDensity\_isRational}, \texttt{countingDensity\_eq\_sum\_coeff},
\texttt{coeff\_isRational}, \texttt{C}, \texttt{undecidedVanishes},
\texttt{nodeMultiplicativity}, \texttt{stratum\_tendsto\_C}, \texttt{boxVolume\_eq},
\texttt{residualBoxCount}, \texttt{residualBoxCount\_eq\_factor},
\texttt{clusterCount\_boxSum}, \texttt{clusterCount\_rec},
\texttt{OMInduction.clusterVol\_isRational} $\to$
$\{\texttt{propext},\ \texttt{Classical.choice},\ \texttt{Quot.sound}\}$ (core only).
\end{itemize}

\subsection{Non-vacuity (the witness, resolved)}\label{sec:vacuity}

An adversarial audit flagged that \texttt{Witness.DensityFoundation\_nonempty}
inhabits only \texttt{DensityFoundation} (density $\equiv0$), which is \emph{necessary
but not sufficient} for \texttt{goal\_theorem\_montes}: its hypothesis
$\texttt{hbridge}$ \emph{couples} $F$ to a \texttt{MontesData}\ $D$, and the trivial
density-$0$ foundation is incompatible with any non-trivial $D$. This is now
\textbf{resolved}: \texttt{Witness.montes\_bundle\_nonempty} (core-only) exhibits a
complete coupled witness ($q=2,n=0$: \texttt{trivM : CountingModel},
\texttt{trivD : MontesData} with the new box-wise fields --- $\texttt{boxMeasure}
:= \texttt{countCellCoeff}$ so \texttt{boxHaarEquidist} is \texttt{rfl} and
\texttt{nodeMeasure\_boxSum} is the constant limit --- and \texttt{trivF} whose
density \emph{is} the count-native tree-sum, so \texttt{hbridge} holds by
\texttt{rfl}), with \emph{non-trivial} density $g_\sigma\equiv1$. The instance
\texttt{Witness.montes\_full\_instance} (= \texttt{goal\_theorem\_montes} applied to
the bundle) is sorry-free with footprint
$\{\texttt{propext},\ \texttt{Classical.choice},\ \texttt{tame\_functionalEquation},\
\texttt{Quot.sound}\}$. So \texttt{goal\_theorem\_montes} is demonstrably non-vacuous
over genuine inhabited types, and the new minimal box-wise axiom is itself
satisfiable. The \texttt{DensityFoundation\_nonempty} docstring was corrected to drop
the prior \texttt{goal\_theorem\_montes} over-claim and point at the coupled witness.

\paragraph{The measure wall stays honest.} This witness crosses \emph{no} measure
wall: \texttt{trivM} is a trivial degree-$0$ model, \emph{not} the paper's $p$-adic
$\rho$. No genuine $p$-adic Haar instance of \texttt{boxMeasure} is (or can be)
constructed in mathlib v4.31.0. The Lean development therefore proves the
\emph{conditional} statement ``\emph{given} the GMN count shadow and the minimal
box-Haar identification, the counting density is a uniform rational, palindromic
function of $q$'' --- a faithful conditional theorem with the assumed ledger above,
not an unconditional construction of $\rho$.

\section*{Verification artifacts}
\begin{itemize}[nosep]
\item \texttt{notes/decomp\_blueprint\_verify.py} --- EXACT (rational-arithmetic)
graded enumeration over $(\mathbb Z/p^N)^s$, three blocks: \textbf{(A)} the $s=2$
single-geometric collapse (constant $a=(q-1)/q$, ratio $q^{-w(2)}=q^{-2}$, limit
$q/(q+1)$, inside the rigorous bracket) at $p\in\{2,3,5\}$; \textbf{(B)} the
\emph{deliberate demonstration} that the single-geometric model fails for $s=3$
($a_N$ strictly increases; increment ratio $\to q^{-2}$, not $q^{-w(3)}$);
\textbf{(C)} the \emph{nested recursion} reproducing $I_2,I_3,I_4$ symbolically and
landing $I_3$ inside the rigorous bracket $[L_N,L_N+U_N]$ ($8/31$ at $p=2$,
$189/484$ at $p=3$). Also verifies the exponent $w(s)=s(s+1)/2-1=2,5,9,14,20$.
\item \texttt{notes/r1\_verify.py} --- the symbolic backbone: the \emph{bottom-up,
non-circular} \texttt{I\_recursion} (base $I_0=I_1=1$) giving
$I_2=q/(q+1)$, $I_3=q^3(q^2-q+1)/((q+1)\Phi_5)$, $I_4$, exponents $w=2,5,9$, and the
\emph{projective} assembly $\rho_{\mathrm{proj}}(3,1^3;q)=(q^2+1)^2/(6\Phi_5)$
(symmetry factor $3!$, divided by $|\mathbb P^3|$ --- \emph{not} the monic density of
Def.~\ref{def:count}, see flag~C3); Vieta Jacobian $=$ Vandermonde, $\disc=V^2$ (the
R1 inputs to Lemma~\ref{lem:rescale}).
\item \texttt{forum-sigma/uniform-rationality/rpoint\_exact.py} --- the exact,
never-Monte-Carlo $r$-point leading-constant enumerator (two-layer exact
$\mu_k$ with a rigorous bracket and self-similar tail extrapolation), the
anti-fooling ground truth.
\end{itemize}

All numeric ``matches'' above are \emph{exact rational equalities} across several
primes (so a bounded-degree rational function of $q$ is pinned), or rigorous
\emph{exact brackets} $L_N\le\text{value}\le L_N+U_N$ that contain the claimed
rational; wild $p=2,3$ included. No floats or Monte-Carlo enter any acceptance test.

\end{document}
